\subsection{Proof of theorem 2}
\label{app:no-vanishing-gradients-proof}

\begin{proof}[Proof of Theorem~\ref{thm:no-vanishing-gradients}]
Fix $x\in\mathcal X$ and $\mu>0$. Fix a recourse matrix $W_f$ and consider the decision-focused loss
\[
L_\mu((W_f,h,q),x)
=
G_\mu(\hat z_\mu^\#(W_f,h,q),x).
\]
We prove that this loss has no spurious stationary points in the learned scenario components $h$ and $q$. The componentwise invexity statement then follows from the standard characterization of continuously differentiable invex functions: a continuously differentiable function is invex precisely when every stationary point is globally optimal.

All gradients below are taken on the relative interior of the synthetic scenario domain where the single-scenario barrier problem is well defined and the derivative formulas from Appendix~\ref{app:closed-form-derivatives} apply.

\paragraph{The rank-condition case.}
Consider the following generalized single-scenario smoothed problem:
\[
\hat z_\mu(h,q)
=
\argmin_z
\left\{
F(z)+Q_\mu(z;W_f,h,q)
\right\},
\]
where $F$ is a differentiable convex first-stage objective on the relative interior of a polytope $Z$, and
\[
Q_\mu(z;W_f,h,q)
=
\min_{y>0}
\left\{
q^\top y-\mu\sum_k\log y_k:
W_fy=h-Tz
\right\}.
\]
Let the corresponding true downstream objective be $G_\mu(z,x)$, assumed convex and differentiable on the same relative interior, with unique optimizer $z_\mu^\#(x)$. Define
\[
L_\mu(h,q,x):=G_\mu(\hat z_\mu(h,q),x).
\]
Assume that $T$ has no nonzero feasible-direction kernel, meaning
\[
d\in\operatorname{lin}(Z-Z)
\quad\text{and}\quad
Td=0
\quad\Longrightarrow\quad
d=0.
\]
Equivalently, every nonzero feasible first-stage movement changes the second-stage right-hand side.

We show that
\[
\nabla_q L_\mu(h,q,x)=0
\quad\text{and}\quad
\nabla_h L_\mu(h,q,x)=0
\]
if and only if
\[
\hat z_\mu(h,q)=z_\mu^\#(x).
\]
In fact, under the rank condition, $\nabla_hL_\mu(h,q,x)=0$ alone already implies $\hat z_\mu(h,q)=z_\mu^\#(x)$.

Let
\[
\hat z=\hat z_\mu(h,q),
\qquad
v=\nabla_zG_\mu(\hat z,x).
\]
By the chain rule,
\[
\nabla_h L_\mu(h,q,x)
=
\left(\frac{\partial\hat z}{\partial h}\right)^\top v.
\]
From the derivative formula established earlier, the $h$-Jacobian has the form
\[
\frac{\partial\hat z}{\partial h}
=
PT^\top S^{-1},
\]
where $P$ is the projected inverse-Hessian operator for the first-stage feasible directions, and $S$ is the positive definite sensitivity matrix arising from the coupled first- and second-stage KKT system. Therefore,
\[
\nabla_h L_\mu(h,q,x)=0
\]
implies
\[
S^{-1}TPv=0.
\]
Since $S$ is invertible,
\[
TPv=0.
\]
Now $Pv$ is a feasible first-stage direction: $P$ maps into $\operatorname{lin}(Z-Z)$. By the rank assumption on $T$, the equality $TPv=0$ implies
\[
Pv=0.
\]
The operator $P$ is positive definite on the feasible-direction space $\operatorname{lin}(Z-Z)$, and its null space is the orthogonal complement of that feasible-direction space. Hence $Pv=0$ means exactly that $v$ has no component along any feasible first-stage direction. Equivalently,
\[
d^\top\nabla_zG_\mu(\hat z,x)=0
\qquad
\text{for every feasible direction }d\in\operatorname{lin}(Z-Z).
\]
Thus $\hat z$ satisfies the first-order optimality condition for minimizing the convex function $G_\mu(\cdot,x)$ over $Z$. Since $\mu>0$ and the smoothed problem has a unique optimizer,
\[
\hat z=z_\mu^\#(x).
\]
This proves the forward direction.

Conversely, suppose
\[
\hat z=z_\mu^\#(x).
\]
Then $\hat z$ minimizes $G_\mu(\cdot,x)$ over $Z$. Hence
\[
d^\top\nabla_zG_\mu(\hat z,x)=0
\qquad
\text{for every feasible direction }d\in\operatorname{lin}(Z-Z).
\]
Both Jacobians $\partial\hat z/\partial h$ and $\partial\hat z/\partial q$ take values in the feasible-direction space $\operatorname{lin}(Z-Z)$, because differentiating the equality constraints defining $Z$ gives $A\,d\hat z=0$. Therefore,
\[
\left(\frac{\partial\hat z}{\partial h}\right)^\top
\nabla_zG_\mu(\hat z,x)=0,
\qquad
\left(\frac{\partial\hat z}{\partial q}\right)^\top
\nabla_zG_\mu(\hat z,x)=0.
\]
By the chain rule,
\[
\nabla_h L_\mu(h,q,x)=0,
\qquad
\nabla_q L_\mu(h,q,x)=0.
\]
Thus, under the rank condition on $T$,
\[
\nabla_q L_\mu(h,q,x)=0
\quad\text{and}\quad
\nabla_h L_\mu(h,q,x)=0
\]
if and only if
\[
\hat z_\mu(h,q)=z_\mu^\#(x).
\]
Finally, since
\[
L_\mu(h,q,x)=G_\mu(\hat z_\mu(h,q),x),
\]
and $G_\mu(\cdot,x)$ is minimized uniquely at $z_\mu^\#(x)$, we also have
\[
\hat z_\mu(h,q)=z_\mu^\#(x)
\]
if and only if
\[
(W_f,h,q)\in\argmin L_\mu((W,h,q),x).
\]
This proves the rank-condition case.

\paragraph{The general case.}
We now remove the rank condition on $T$ by projecting away the part of $z$ that does not affect the second stage.

Let
\[
Z=\{z\geq 0:\;Az=b\},
\qquad
Z^\circ=Z\cap\R^{n_1}_{++}
\]
denote the barrier domain. Let
\[
V=\operatorname{lin}(Z-Z)
\]
be the feasible-direction space of the first-stage polytope. The obstruction to the rank condition is the subspace
\[
K=V\cap\ker T.
\]
These are precisely the feasible first-stage directions along which $z$ can move without changing the second-stage right-hand side $h-Tz$.

We quotient out these directions. Equivalently, identify two feasible first-stage points $z_1,z_2\in Z$ whenever
\[
Tz_1=Tz_2.
\]
Let
\[
u=Tz
\]
be the reduced first-stage variable, and define the reduced feasible set
\[
Z_{\mathrm{red}}:=T(Z).
\]
Since $Z$ is a polytope and $T$ is linear, $Z_{\mathrm{red}}$ is also a polytope. Its relative interior is the image of the relative interior of $Z$ under $T$.

For $u\in Z_{\mathrm{red}}$, define the fiber
\[
Z(u):=\{z\in Z:\;Tz=u\}.
\]
Now define the reduced first-stage cost by partial minimization over the fiber:
\[
F_\mu(u)
:=
\min_{z\in Z(u)\cap\R^{n_1}_{++}}
\left\{
c^\top z-\mu\sum_j\log z_j
\right\}.
\]
The function $F_\mu$ is convex. Indeed, it is the value function obtained by partially minimizing a jointly convex function over affine slices. More explicitly, if $u_1,u_2\in Z_{\mathrm{red}}$ and $\alpha\in[0,1]$, then feasible fiber points over $u_1$ and $u_2$ convexly combine to a feasible fiber point over $\alpha u_1+(1-\alpha)u_2$, and convexity of the objective gives
\[
F_\mu(\alpha u_1+(1-\alpha)u_2)
\leq
\alpha F_\mu(u_1)+(1-\alpha)F_\mu(u_2).
\]
Because $\mu>0$, the fiber minimizer is unique whenever the fiber has nonempty barrier-relative interior. Denote this unique fiber minimizer by
\[
\phi_\mu(u)
:=
\argmin_{z\in Z(u)\cap\R^{n_1}_{++}}
\left\{
c^\top z-\mu\sum_j\log z_j
\right\}.
\]

Now observe that the second-stage value depends on $z$ only through $Tz$. Therefore, for every fixed scenario $(W,h,q)$,
\[
Q_\mu(z;W,h,q)=Q_\mu^{\mathrm{red}}(u;W,h,q),
\qquad
u=Tz,
\]
where
\[
Q_\mu^{\mathrm{red}}(u;W,h,q)
=
\min_{y>0}
\left\{
q^\top y-\mu\sum_k\log y_k:
Wy=h-u
\right\}.
\]
Hence the original single-scenario optimization problem
\[
\min_{z\in Z^\circ}
\left\{
c^\top z-\mu\sum_j\log z_j+Q_\mu(z;W_f,h,q)
\right\}
\]
is equivalent to the reduced problem
\[
\min_{u\in Z_{\mathrm{red}}}
\left\{
F_\mu(u)+Q_\mu^{\mathrm{red}}(u;W_f,h,q)
\right\}.
\]
If $\hat u_\mu(h,q)$ denotes the unique optimizer of the reduced problem, then the original optimizer is recovered by
\[
\hat z_\mu(W_f,h,q)=\phi_\mu(\hat u_\mu(h,q)).
\]

The same reduction applies to the true stochastic objective. Define
\[
G_\mu^{\mathrm{red}}(u,x)
:=
F_\mu(u)
+
\E_{\boldsymbol\xi\sim P(\cdot\mid x)}
\left[
Q_\mu^{\mathrm{red}}(u;W(\boldsymbol\xi),h(\boldsymbol\xi),q(\boldsymbol\xi))
\right].
\]
Then, for every $z\in Z^\circ$ with $u=Tz$,
\[
G_\mu(z,x)\geq G_\mu^{\mathrm{red}}(u,x),
\]
with equality precisely when $z=\phi_\mu(u)$. Therefore the original stochastic optimizer satisfies
\[
z_\mu^\#(x)=\phi_\mu(u_\mu^\#(x)),
\]
where
\[
u_\mu^\#(x)
=
\argmin_{u\in Z_{\mathrm{red}}}G_\mu^{\mathrm{red}}(u,x).
\]
Likewise, the original decision-focused loss can be written entirely in reduced coordinates:
\[
\begin{aligned}
L_\mu((W_f,h,q),x)
&=G_\mu(\hat z_\mu(W_f,h,q),x)\\
&=G_\mu^{\mathrm{red}}(\hat u_\mu(h,q),x).
\end{aligned}
\]
Thus the original scenario-parametrized loss is exactly the reduced scenario-parametrized loss, composed with the fiber minimizer $\phi_\mu$.

Now consider the reduced problem. Its decision variable is $u$, and the second-stage right-hand side is $h-u$. Therefore the reduced technology map is the identity map on the reduced feasible-direction space. In particular, it satisfies the rank condition from the rank-condition case automatically: there is no nonzero reduced direction $du$ with $du=0$.

The reduced problem also has a convex first-stage cost $F_\mu$, exactly of the type allowed above. Therefore, applying the rank-condition case to the reduced problem gives
\[
\nabla_q L_\mu((W_f,h,q),x)=0
\quad\text{and}\quad
\nabla_h L_\mu((W_f,h,q),x)=0
\]
if and only if
\[
\hat u_\mu(h,q)=u_\mu^\#(x).
\]
Using the recovery map $\phi_\mu$, this is equivalent to
\[
\hat z_\mu(W_f,h,q)
=
\phi_\mu(\hat u_\mu(h,q))
=
\phi_\mu(u_\mu^\#(x))
=
z_\mu^\#(x).
\]
Thus
\[
\nabla_q L_\mu((W_f,h,q),x)=0
\quad\text{and}\quad
\nabla_h L_\mu((W_f,h,q),x)=0
\]
if and only if
\[
\hat z_\mu(W_f,h,q)=z_\mu^\#(x).
\]

It remains to identify the global minimizers over the synthetic scenario domain. Since
\[
L_\mu((W_f,h,q),x)
=
G_\mu(\hat z_\mu(W_f,h,q),x),
\]
and $G_\mu(\cdot,x)$ has unique minimizer $z_\mu^\#(x)$, every scenario satisfying
\[
\hat z_\mu(W_f,h,q)=z_\mu^\#(x)
\]
is globally optimal for the scenario-parametrized loss.

Conversely, by Theorem~\ref{thm:realizability}, there exists at least one synthetic scenario with recourse matrix $W_f$ whose induced decision is $z_\mu^\#(x)$. Hence the minimum achievable scenario-parametrized loss is exactly
\[
G_\mu(z_\mu^\#(x),x).
\]
Therefore, if $(W_f,h,q)$ is a global minimizer of $L_\mu$, then
\[
G_\mu(\hat z_\mu(W_f,h,q),x)
=
G_\mu(z_\mu^\#(x),x).
\]
By uniqueness of $z_\mu^\#(x)$, this implies
\[
\hat z_\mu(W_f,h,q)=z_\mu^\#(x).
\]
Combining the above implications gives the equivalence
\[
\nabla_q L_\mu((W_f,h,q),x)=0
\quad\text{and}\quad
\nabla_h L_\mu((W_f,h,q),x)=0
\]
if and only if
\[
(W_f,h,q)\in\argmin L_\mu((W,h,q),x)
\]
if and only if
\[
\hat z_\mu(W_f,h,q)=z_\mu^\#(x).
\]
This is exactly the claim of Theorem~\ref{thm:no-vanishing-gradients}.

Since $L_\mu$ is continuously differentiable in $q$ and $h$ on the relevant synthetic domain, and every stationary point is globally optimal, the standard first-order characterization of invexity gives componentwise invexity in the learned scenario components $q$ and $h$.
\end{proof}
